\chapter{Intermediate Extension and Intersection Complexes}
\label{chap:intermediate-extension}

An extension by zero has too many quotients on the boundary; a direct image
has too many boundary subobjects.  The intermediate extension is the unique
extension with neither defect.  This minimality is what makes it compatible
with duality and what singles out intersection cohomology.

Fix an open immersion \(j:U\hookrightarrow X\), with complementary closed
immersion \(i:Z\hookrightarrow X\).

\section{The image in the perverse heart}

For \(\cP\in\Perv(U)\), adjunction gives a canonical morphism
\[
  {}^pj_!\cP=\pH^0(j_!\cP)
  \longrightarrow
  {}^pj_*\cP=\pH^0(Rj_*\cP).
\tag{9.1}\label{eq:standard-to-costandard}
\]

\begin{definition}[Intermediate extension]
\label{def:intermediate-extension}
The intermediate extension of \(\cP\) is the image in the abelian category
\(\Perv(X)\):
\[
  j_{!*}\cP
  =
  \im_{\Perv(X)}
  \bigl({}^pj_!\cP\longrightarrow{}^pj_*\cP\bigr).
\tag{9.2}\label{eq:intermediate-extension}
\]
\end{definition}

\begin{warning}
The symbol “image” in \eqref{eq:intermediate-extension} is meaningful because
\(\Perv(X)\) is abelian.  A triangulated category has no functorial kernels,
cokernels, or images.  Thus \(j_{!*}\) cannot be defined as an image directly
inside \(\Dbc(X)\).
\end{warning}

Restriction to \(U\) is exact and sends both sides of
\eqref{eq:standard-to-costandard} to \(\cP\).  Consequently
\[
  j^*j_{!*}\cP\simeq\cP.
\tag{9.3}\label{eq:jstar-jmiddle}
\]

\begin{theorem}[Minimal-extension characterization]
\label{thm:minimal-extension-characterization}
Let \(\cQ\in\Perv(X)\) and suppose \(j^*\cQ\simeq\cP\).  Then
\(\cQ\simeq j_{!*}\cP\) if and only if
\begin{enumerate}
\item \(\cQ\) has no nonzero quotient supported on \(Z\);
\item \(\cQ\) has no nonzero subobject supported on \(Z\).
\end{enumerate}
Equivalently,
\[
  {}^pi^*\cQ=0,\qquad {}^pi^!\cQ=0.
\tag{9.4}\label{eq:middle-boundary-vanishing}
\]
The isomorphism extending the identity of \(\cP\) is unique.
\end{theorem}

\begin{proof}
The map \({}^pj_!\cP\to\cQ\) is forced by adjunction.  Since its restriction
to \(U\) is an isomorphism, its cokernel is supported on \(Z\); condition (1)
makes it surjective.  Dually, the adjunction map
\(\cQ\to{}^pj_*\cP\) has kernel supported on \(Z\); condition (2) makes it
injective.  Their composite is the canonical map
\eqref{eq:standard-to-costandard}, so \(\cQ\) identifies with its image.

Conversely, any quotient of that image supported on \(Z\) would be a quotient
of \({}^pj_!\cP\) whose restriction to \(U\) is zero.  Adjunction shows such a
quotient is zero.  The subobject assertion follows dually from
\({}^pj_*\cP\).  Recollement identifies maximal boundary quotients and
subobjects with \(i_*{}^pi^*\cQ\) and \(i_*{}^pi^!\cQ\), proving
\eqref{eq:middle-boundary-vanishing}.  The same argument applied to the graph
of two extensions gives uniqueness.
\end{proof}

\begin{corollary}[Transitivity]
\label{cor:intermediate-extension-transitivity}
If \(V\xrightarrow{k}U\xrightarrow{j}X\) are open immersions, then
\[
  (j\circ k)_{!*}\simeq j_{!*}k_{!*}.
\]
\end{corollary}

\begin{proof}
Both sides restrict to the original perverse sheaf on \(V\), and neither has a
subobject or quotient supported on \(X\setminus V\).  Apply
\cref{thm:minimal-extension-characterization}.
\end{proof}

\begin{warning}
The functor \(j_{!*}\) preserves monomorphisms and epimorphisms: a kernel or
cokernel created at the boundary would contradict the minimal-extension
property.  It is nevertheless not exact in general, because middle exactness
of a short exact sequence can fail.  What is canonical is the image
construction in the perverse heart.
\end{warning}

\section{Compatibility with Verdier duality}

\begin{proposition}[Self-duality of intermediate extension]
\label{prop:duality-intermediate-extension}
There is a canonical isomorphism
\[
  \D_X(j_{!*}\cP)\simeq j_{!*}(\D_U\cP).
\tag{9.5}\label{eq:duality-jmiddle}
\]
\end{proposition}

\begin{proof}
Verdier duality is exact on the perverse heart and exchanges
\({}^pj_!\) with \({}^pj_*\).  It therefore sends the image of
\({}^pj_!\cP\to{}^pj_*\cP\) to the coimage of the dual map.  In an abelian
category image and coimage agree.  Alternatively, duality interchanges
boundary subobjects and boundary quotients, so the result follows immediately
from \cref{thm:minimal-extension-characterization}.
\end{proof}

\section{Finite length and simple perverse sheaves}

\begin{theorem}[BBD finite-length theorem]
\label{thm:perverse-finite-length}
The category \(\Perv(X,\Qbarl)\) is both Noetherian and Artinian.  Every
perverse sheaf has finite length.
\end{theorem}

\begin{proof}
Induct on a stratification of \(X\).  On a smooth connected stratum, a
perverse sheaf is generically a shifted lisse sheaf, and its subobjects give
subrepresentations of a finite-dimensional continuous representation; these
satisfy both chain conditions.  Shrink to make the generic constituent lisse.
The kernel and cokernel of the resulting standard or costandard extension are
supported on the smaller closed complement, where induction applies.  The
recollement exact sequences then prevent an infinite ascending or descending
chain on \(X\).  This is the devissage in BBD 4.3.1.
\end{proof}

\begin{theorem}[Classification of simple objects]
\label{thm:simple-perverse-classification}
Let \(S\subset X\) be a smooth connected locally closed subvariety of
dimension \(d\), let \(\cL\) be an irreducible lisse \(\Qbarl\)-sheaf on
\(S\), and let \(j:S\hookrightarrow\overline S\) and
\(i:\overline S\hookrightarrow X\).  Then
\[
  i_*j_{!*}(\cL[d])
\tag{9.6}\label{eq:simple-IC-object}
\]
is a simple perverse sheaf on \(X\).  Every simple perverse sheaf is obtained
uniquely in this way, up to shrinking \(S\) inside a dense open of its closure
and identifying the resulting local systems.
\end{theorem}

\begin{proof}
Any nonzero subobject or quotient of \(j_{!*}\cL[d]\) restricts to a nonzero
subobject or quotient of \(\cL[d]\).  Irreducibility makes that restriction
all of \(\cL[d]\), and the kernel or cokernel is then supported on the
boundary; minimality makes it zero.  Thus \eqref{eq:simple-IC-object} is
simple.

Conversely, let \(\cP\) be simple and choose a maximal-dimensional stratum
\(S\) on which it is nonzero and lisse after the dimension shift.  Simplicity
forces the generic local system to be irreducible.  The adjunction morphisms
factor \(\cP\) between the standard and costandard extensions, hence produce a
nonzero map \(j_{!*}\cL[d]\to\cP\).  A nonzero map between simple objects is
an isomorphism.  The support and generic local system recover
\(\overline S\) and \(\cL\).
\end{proof}

\section{Intersection complexes}

\begin{definition}[Intersection complex]
\label{def:intersection-complex}
Let \(X\) be irreducible of dimension \(d\), let
\(j:U\hookrightarrow X\) be a smooth dense open, and let \(\cL\) be lisse on
\(U\).  The perverse-normalized intersection complex is
\[
  \IC_X(\cL)=j_{!*}(\cL[d]).
\tag{9.7}\label{eq:IC-definition}
\]
For the constant local system write \(\IC_X=\IC_X(\Qbarl)\).
\end{definition}

Transitivity shows that \(\IC_X(\cL)\) is independent of replacing \(U\) by a
smaller dense smooth open.  When \(\cL\) is irreducible it is a simple
perverse sheaf.  If \(X\) is reducible, define the intersection complex as
the direct sum over irreducible components, with the desired coefficient
local systems on dense smooth opens.

\begin{proposition}[Smooth normalization and duality]
\label{prop:IC-smooth-duality}
\begin{enumerate}
\item If \(X\) is smooth of pure dimension \(d\) and \(\cL\) is lisse on
      all of \(X\), then
      \[
        \IC_X(\cL)=\cL[d].
      \tag{9.8}\label{eq:IC-smooth}
      \]
\item For \(U\) smooth of dimension \(d\),
      \[
        \D_X\IC_X(\cL)
        \simeq \IC_X(\cL^\vee(d)).
      \tag{9.9}\label{eq:IC-duality}
      \]
\end{enumerate}
\end{proposition}

\begin{proof}
For the first statement take \(j=\id_X\).  For the second, smooth Poincare
duality gives
\[
  \D_U(\cL[d])\simeq\cL^\vee(d)[d].
\]
Apply \cref{prop:duality-intermediate-extension}.
\end{proof}

There are two common normalizations.  In these notes \(\IC_X(\cL)\) is
perverse, so on a smooth \(d\)-fold it is \(\cL[d]\).  Some authors shift back
and call \(\IC_X(\cL)[-d]\) the intersection complex.  Every weight formula
below uses the perverse normalization \eqref{eq:IC-definition}.

\section{The curve model}

Let \(C\) be a smooth connected curve and \(j:U\hookrightarrow C\) a dense
open.  For a lisse sheaf \(\cL\) on \(U\),
\[
  \IC_C(\cL)=j_{!*}(\cL[1])\simeq (j_*\cL)[1],
\tag{9.10}\label{eq:curve-IC}
\]
where \(j_*\cL\) on the right is the ordinary sheaf direct image.  At a
boundary point \(x\), its stalk is the inertia-invariant space
\((\cL_{\bar\eta})^{I_x}\).  Formula \eqref{eq:curve-IC} follows by checking
the support and cosupport conditions and the absence of punctual subobjects
and quotients.

Consequently
\[
  \mathbb H^r(C,\IC_C(\cL))
  =H^{r+1}(C,j_*\cL).
\tag{9.11}\label{eq:curve-IH}
\]
The degree-zero intersection cohomology
\(\mathbb H^0(C,\IC_C(\cL))=H^1(C,j_*\cL)\) is the parabolic cohomology that
appears in the curve weight theorem.

\begin{example}[A nodal curve]
Let \(X\) be an irreducible nodal proper curve with normalization
\(\nu:\widetilde X\to X\).  Since \(\nu\) is finite,
\[
  \IC_X\simeq\nu_*\Qbarl[1].
\]
At the node \(x\),
\[
 \cH^{-1}(\nu_*\Qbarl[1])_x\simeq\Qbarl^2,
 \qquad
 \cH^{-1}\bigl((\Qbarl)_X[1]\bigr)_x\simeq\Qbarl.
\]
Thus the shifted constant complex identifies the two branches through the
diagonal gluing condition, whereas the intersection complex retains the two
normalization branches separately.  The latter is the self-dual pure object
appropriate to the singular curve.
\end{example}

\section{Where weights will enter}

Nothing in this chapter proves that \(\IC_X(\cL)\) is pure.  Once weights are
constructed, the correct shift will be
\[
  \cL\text{ pure of weight }w
  \quad\Longrightarrow\quad
  \IC_X(\cL)\text{ pure of weight }w+d.
\tag{9.12}\label{eq:IC-weight-preview}
\]
The dimension term \(d\) is not cosmetic: it is forced by the convention that
ordinary cohomology sheaves of a pure complex of weight \(v\) have weights
\(v+a\) in degree \(a\).  The proof of \eqref{eq:IC-weight-preview} is delayed
until \cref{chap:perverse-weights}; proper cohomological purity is delayed
further, until the Deligne--Laumon direct-image theorem has been proved.

\begin{formalization}
The minimum reusable API consists of:
\begin{enumerate}
\item heart-level standard and costandard extensions;
\item the canonical natural transformation
      \({}^pj_!\Rightarrow{}^pj_*\);
\item \(j_{!*}\) as its abelian image;
\item the equivalent no-boundary-subobject/no-boundary-quotient predicate;
\item transitivity and Verdier compatibility;
\item finite length and the classification of simple objects;
\item \(\IC_X(\cL)\) with its smooth and curve normalization tests.
\end{enumerate}
Choice independence should be an isomorphism characterized by the universal
property, not a definitional equality tied to a selected stratification.
\end{formalization}
